\cotefile{57}
\scanpage{20}
Dear Murre,\\

I am glad to hear that you are still willing to \rlap{\deleted{\unsure{make}}}gi\unsure{v}e the talk on unramified functors. Here what I can say to your questions.

1. The theorem about passage to quotient I alluded to is the foll\unsure{owing|hwing}:

\uline{Theorem} Let $f:X\handwritten{\to}Y$ be a morphisme of $S$-preschemes, assume either $X$ and $Y$ \tikz[baseline=(s.base)]{\node[inner sep=0pt](s){of};\node[anchor=south west,inner sep=1pt](a) at ([xshift=-.75em,yshift=6pt]s.north west){\inserted{locally}};\draw ([xshift=-3pt,yshift=-3pt]s.south west) to[out=100,in=180] ([xshift=0pt,yshift=-2pt]a.south west)--([xshift=0pt,yshift=-2pt]a.south east) to[out=0,in=-70] ([xshift=5pt,yshift=5pt]a.south east);} finite presentation over $\rlap{\deleted{$Y$}}\handwritten{S}$, or $Y$ loc noeth and $X$ \tikz[baseline=(s.base)]{\node[inner sep=0pt](s){of};\node[anchor=south west,inner sep=1pt](a) at ([xshift=-2em,yshift=6pt]s.north west){\inserted{locally}};} finite type over $Y\handwritten{\unsure{'}}$. Assume t\rlap{a}h\rlap{h}at the equivalence relation $\mathord{\tikz[baseline=(s.base)]{\node[inner sep=0pt](s){$X$};\node[anchor=south west,inner sep=1pt](a) at ([xshift=-1em,yshift=6pt]s.north west){\inserted{$R=$}};}}\mathbin{x}_{Y}X$ define\rlap{\deleted{e}}\handwritten{d} by $f$ is flat over $X$ i.e. $\mathrm{pr}_{1}:X\mathbin{x}_{Y}X\handwritten{\to}X$ is flat. Then the quotient $X/R$ exis\unsure{t}% clipped by right edge; no final s is visible
 in the strongest reasonable sense, \tikz[baseline=(s.base)]{\node[inner sep=0pt](s){i};\draw[line width=.6pt] ([xshift=-1.5pt,yshift=3pt]s.north)--([xshift=1.5pt,yshift=3pt]s.north) ([yshift=3pt]s.north)--([yshift=4.2pt]s.north);\fill ([yshift=-2pt]s.south) circle (.5pt);\draw[line width=.7pt] ([xshift=-1pt,yshift=-3.2pt]s.south)--([xshift=.4pt,yshift=-3.5pt]s.south);}.e. one can factor $f$ into a compositum $X\quad Z\quad Y$, with $\unsure{\mathbb{X}|\mathbb{Y}}\quad Z$ faithfully flat locally of finite presentation, $Z$ locally of finite \tikz[baseline=(s.base)]{\node[inner sep=0pt](s){\deleted{type}};\node[anchor=south west,inner sep=1pt](a) at ([xshift=0em,yshift=6pt]s.north west){\inserted{pres.}};} over $\mathord{\rlap{$Y$}\raisebox{-.15ex}{$S$}}$ (in fact of finite presentation over $S$ if $X$ is so) and $Z\handwritten{\mathrel{\tikz[baseline=(s.base)]{\node[inner sep=0pt](s){$\to$};\draw[gray,line width=.25pt] ([xshift=1pt,yshift=.5pt]s.west) to[out=30,in=180] ([xshift=5pt,yshift=2pt]s.west);}}}Y$ a monomorphism.

Of course the factorization is unique, andthe theorem can be expres-
sed by saying that the quotient sheaf (for the fpqc topology) $X/R$ is representable. T\tikz[baseline=(s.base)]{\node[inner sep=0pt](s){h};\draw[line width=.7pt] ([xshift=.2pt,yshift=2pt]s.south west) to[out=150,in=130] ([xshift=.6pt,yshift=-5pt]s.south west) to[out=70,in=140] ([xshift=-1pt,yshift=-2pt]s.south east)--([xshift=-.2pt,yshift=-5.5pt]s.south east); }at is in fact how the theorem is proved.

Raynaud has recently made a very nice (and non trivial) application of this theorem, by proving the following: if $S$ is the spectrum of a discrete valuation ring, $G$ a group prescheme of finite type over $\mathord{\tikz[baseline=(s.base)]{\node[inner sep=0pt](s){$S$};\fill ([xshift=1pt,yshift=-3pt]s.south) circle (.25pt);\fill ([xshift=-1pt,yshift=-5pt]s.south) circle (.22pt);}}$, $H$ a closed \tikz[baseline=(s.base)]{\node[inner sep=0pt](s){sub-group};\node[anchor=south west,inner sep=1pt](a) at ([xshift=-2em,yshift=6pt]s.north west){\inserted{and flat}};} scheme, such that $G_{t}/H_{t}$ is \deleted{affi} quasi-affine (where $t$ is the generic point of $S$) then $G/H$ is representable as a quasi-affine and flat $S$-scheme, which is even affine if $H$ is invariant (i.e. if $G$ is a flat group scheme of finite type with affine generic fibre, th\rlap{e}an $G$ is affine). This extends immediately to a b\kern-.2em\raisebox{-.6ex}{\inserted{\handwritten{\unsure{a}}}}\kern-.2em se which is regular of dim one. Raynaud is now trying to extend his construction to the c\rlap{q}ase when he drops the quasi-affinness assumption, namely to construct still $G/H$
% archive sheet number: 19
